El Fun-Fly-Stick és un enginyós generador d'electricitat estàtica que funciona amb piles i que permet fer levitar objectes metàl·lics lleugers. Quan es prem el botó es genera una càrrega estàtica negativa a la vareta. Tocant l'objecte amb la vareta, part d'aquesta càrrega es transfereix a l'objecte de manera que tots dos queden carregats i l'objecte levita per sobre la vareta. (En fer els càlculs, considereu la vareta i l'objecte metàl·lic com si fossin càrregues puntuals.)

\figura[0.25]{fig1.png}

\begin{apartat}{1,25}
Suposem que, després de tocar l'objecte, la vareta i l'objecte tenen la mateixa càrrega. Feu un dibuix de la situació de la vareta i l'objecte metàl·lic levitant i dibuixeu les forces sobre l'objecte. Calculeu quina és la càrrega electroestàtica de l'objecte metàl·lic si aquest té una massa de 10~g i es troba a una distància de 55~cm de la vareta.
\end{apartat}

\begin{apartat}{1,25}
Considereu ara una altra situació, en què tant la vareta com l'objecte metàl·lic estan carregats amb una càrrega $q = -2\ \mathrm{\mu C}$ i separats a una distància de 60~cm. Calculeu el mòdul del camp elèctric al punt central de la línia que els uneix i el potencial elèctric en aquest punt. Determineu quin és el treball que haurà de fer una força externa per a portar un electró des de l'infinit fins a aquest punt central.
\end{apartat}

\dades{%
  Gravetat a la superfície de la Terra, $g = 9,8\ \mathrm{m/s^2}$.\\
  $|e| = 1,602\times 10^{-19}\ \mathrm{C}$.\\
  $k = \dfrac{1}{4\pi\varepsilon_0} = 8,99\times 10^{9}\ \mathrm{N\,m^2\,C^{-2}}$.}
